\documentclass[12pt]{article}
\usepackage[brazil]{babel}
\usepackage[utf8]{inputenc}
\usepackage[T1]{fontenc}
\usepackage{amsmath, amssymb}
\usepackage{geometry}
\usepackage{graphicx}
\usepackage{xcolor}

\definecolor{maincolor}{RGB}{0,100,50}
\usepackage{fancyhdr}
\definecolor{forestgreen}{RGB}{0,100,50}
\usepackage[most]{tcolorbox}

\newtcolorbox{block}[1]{%
  colback=green!5,
  colframe=maincolor,
  coltitle=white,
  title=#1,
  fonttitle=\bfseries,
  boxrule=1.5pt,
  arc=3mm,
  left=6pt,
  right=6pt,
  top=6pt,
  bottom=6pt
}

\geometry{margin=2.5cm}
\pagestyle{fancy}
\fancyhf{}

\renewcommand{\headrulewidth}{2pt}
\renewcommand{\footrulewidth}{2pt}

\renewcommand{\headrule}{\hbox to\headwidth{%
  \color{forestgreen}\leaders\hrule height \headrulewidth\hfill}}

\renewcommand{\footrule}{\hbox to\headwidth{%
  \color{forestgreen}\leaders\hrule height \footrulewidth\hfill}}

\setlength{\headheight}{32pt}

\fancyhead[L]{%
  \includegraphics[height=1.2cm]{Universidade_de_Rio_Verde.jpg}
}

\fancyhead[R]{%
  \raggedleft
  \textcolor{forestgreen}{%
    \small
    \textbf{Lista de Exercícios — Provas com Conjuntos}\\
    Me.~Sergio Souza Novak\\
    \textit{Lógica para Engenharia de Software}
  }
}

\fancyfoot[C]{%
  \textcolor{forestgreen}{\thepage}
}

\begin{document}

\vspace*{0.2cm}

\section*{Instruções}

\begin{block}{Utilize definições formais de conjuntos}
Utilize rigorosamente as definições:

\[
A \subseteq B \iff \forall x (x \in A \Rightarrow x \in B)
\]

\[
x \in A \cap B \iff (x \in A \land x \in B)
\]

\[
x \in A \cup B \iff (x \in A \lor x \in B)
\]

e o fato fundamental:

\[
\forall x,\; x \notin \varnothing.
\]
\end{block}

Justifique formalmente todas as demonstrações.

\section*{Exercícios}

\begin{enumerate}

\item Prove que não existe elemento pertencente ao conjunto vazio.

\item Demonstre que não existe conjunto $A$ tal que
\[
A \subseteq \varnothing
\]
e $A$ possua elementos.

\item Prove que, se
\[
x \in A \cap \varnothing,
\]
então ocorre uma contradição.

\item Demonstre que, para qualquer conjunto $A$,
não existe elemento comum entre $A$ e o conjunto vazio.

\item Prove que toda proposição da forma
\[
x \in \varnothing \Rightarrow P(x)
\]
é verdadeira para qualquer proposição $P(x)$.

\item Suponha que $A \subseteq \varnothing$. Prove que $A=\varnothing$.

\item Se $A \cap B=\varnothing$, prove que
\[
x \in A \Rightarrow x \notin B.
\]

\item Prove que
\[
(A \cap \varnothing)\subseteq B
\]
para qualquer conjunto $B$.

\item Demonstre que
\[
A \cup (A \cap \varnothing)=A.
\]

\item Prove a unicidade do conjunto vazio.

\end{enumerate}

\end{document}